% §1151 `scan_math` lists what a math *field* may be, and `\delimiter` is one of
% them, beside a character, `\char`, `\mathchar` and a `\mathchardef`:
%
%     delim_num: begin scan_twenty_seven_bit_int;
%       cur_val:=cur_val div @'10000; goto found;
%       end;
%
% Its 27 bits are class, small variant and large variant. Used as an ordinary
% symbol only the small one applies, which is the value shifted down twelve
% bits — the same arithmetic `main_control` does for a `\delimiter` met directly
% in a formula, and the reason `\mathopen\lVert` sets one glyph rather than a
% grown fence.
%
% We had the arithmetic in `main_control` and not in `scan_math`, so a
% `\delimiter` handed to `\mathopen` and friends fell through to §1151's
% subformula branch, which backs the token up and calls `scan_left_brace`:
%
%     ! Missing { inserted.
%     <recently read> \delimiter
%
% mathtools' `\DeclarePairedDelimiter` puts one there every time. Its unstarred,
% unscaled wrapper is
%
%     \mathopen##1##2\mathclose##3
%
% so `\norm{Z}` from `\DeclarePairedDelimiter{\norm}{\lVert}{\rVert}` comes out
% as `\mathopen\lVert Z\mathclose\rVert`, and `\lVert` is `\delimiter"026B30D`.
% 2608.15840v1 uses `\norm` throughout and stopped at the hundred-error limit
% after four of its thirty-six pages.
\catcode`\{=1 \catcode`\}=2 \catcode`\$=3 \catcode`\^=7 \catcode`\_=8
\font\r=cmr10 \r \hsize=300pt \parindent=0pt
\font\mi=cmmi10 \font\sy=cmsy10 \font\ex=cmex10
\textfont0=\r \scriptfont0=\r \scriptscriptfont0=\r
\textfont1=\mi \scriptfont1=\mi \scriptscriptfont1=\mi
\textfont2=\sy \scriptfont2=\sy \scriptscriptfont2=\sy
\textfont3=\ex \scriptfont3=\ex \scriptscriptfont3=\ex
\showboxdepth=6 \showboxbreadth=20 \tracingonline=0

% The shape mathtools produces: an open atom whose nucleus is a `\delimiter`,
% the body, then a close atom the same way. `"026B30D` is `\lVert`, `"026B30E`
% is `\rVert` — class 0, family 2, character `"0D`/`"0E` once the low twelve
% bits are dropped.
\setbox0=\hbox{$\mathopen\delimiter"026B30D Z\mathclose\delimiter"026B30E$}
\showbox0

% The same delimiters used the ordinary way, straight in the formula rather
% than as a field: `main_control` already took the small variant, and the two
% routes must agree.
\setbox0=\hbox{$\delimiter"026B30D Z\delimiter"026B30E$}
\showbox0

% Every one of §1156's class-changing primitives reads its field the same way.
\setbox0=\hbox{$\mathord\delimiter"026B30D$}
\showbox0
\setbox0=\hbox{$\mathbin\delimiter"026B30D$}
\showbox0
\setbox0=\hbox{$\mathrel\delimiter"026B30D$}
\showbox0

% §1151's other callers: the argument of `\underline`, and of `\radical`.
\setbox0=\hbox{$\underline\delimiter"026B30D$}
\showbox0
\setbox0=\hbox{$\radical"270370 \delimiter"026B30D$}
\showbox0

% A script field, which reads a field too.
\setbox0=\hbox{$x^\delimiter"026B30D$}
\showbox0

% A `\delimiter` whose class bits are not zero — `"426B30D` is class 4, open.
% As a *field* the class is dropped along with the large variant, so this must
% set the same atom as the class-0 one above: §1151 divides the whole 27-bit
% value, class included, by `'10000`.
\setbox0=\hbox{$\mathord\delimiter"426B30D$}
\showbox0

% The control: a braced subformula still goes the other way, through
% `scan_left_brace`. A fix that swallowed everything as a field would lose the
% group.
\setbox0=\hbox{$\mathopen{ab}\mathclose{cd}$}
\showbox0
\end
